\begin{abstract}

In this paper, we make the case that a scientific theory of deep learning is emerging.
By this we mean a theory which characterizes important properties and statistics of the training process, hidden representations, final weights, and performance of neural networks.
We pull together major strands of ongoing research in deep learning theory and identify five growing bodies of work that point toward such a theory:
\begin{enumerate}
    \item solvable idealized settings that provide intuition for learning dynamics in realistic systems;
    \item  tractable limits that reveal insights into fundamental learning phenomena;
    \item simple mathematical laws that capture important macroscopic observables;
    \item theories of hyperparameters that disentangle them from the rest of the training process, leaving simpler systems behind; and
    \item universal behaviors shared across systems and settings which clarify which phenomena call for explanation.
\end{enumerate}

Taken together, these bodies of work share certain broad traits: they are concerned with the dynamics of the training process; they primarily seek to describe coarse aggregate statistics; and they emphasize falsifiable quantitative predictions.
We argue that the emerging theory is best thought of as a \textit{mechanics} of the learning process, and suggest the name \textit{learning mechanics}.
We assert that learning mechanics should be a mathematical theory, grounded in first-principles calculations that closely predict empirics, reliant on well-tested approximations and assumptions, aiming for broad impact across the machine learning stack once it reaches maturity.

We discuss the relationship between this mechanics perspective and other approaches for building a theory of deep learning, including the statistical and information-theoretic perspectives.
In particular, we anticipate a symbiotic and mutually supportive relationship between learning mechanics and the developing discipline of mechanistic interpretability.
Where mechanistic interpretability aims to be the biology of deep learning, learning mechanics should aspire to be its physics, mirroring the complementary relationship between biology and physics in the natural sciences.

We also review and address common arguments that fundamental theory will not be possible or is not important.
We conclude with a portrait of important open directions in learning mechanics and advice for beginners.
We host further introductory materials, perspectives, and open questions at \websiteurl.

















\end{abstract}


\numberlessfootnote{$^*$Correspondence to \texttt{james.simon@berkeley.edu}.}
